% Pista ana-24s-q3 · Anàlisi
% Procedència: PAU Catalunya 2024 · Convocatòria de setembre · Sèrie 3 · Qüestió 3
\documentclass[11pt,a4paper]{article}
\usepackage{pista}
\pistainfo{Catalunya 2024}{Convocatòria de setembre}{Sèrie 3}

\begin{document}

\section*{\textcolor{blauPista}{Qüestió 3}}

\noindent La classe de l'Èlia ha dissenyat el logotip següent per a pintar-lo a la paret de l'institut:

\noindent\textit{(A la figura: $B$ és a l'eix $OY$, i $A$ i $D$ són a la part positiva de l'eix $OX$. La zona puntejada és la compresa entre les dues corbes, des de l'eix $OY$ fins a $C$; la resta del logotip, fins a $D$, és la zona ratllada.)}

\noindent La corba que passa pel punt $A$ és $y = f(x)$, amb $f(x) = x^{3} - 4x^{2} + 4x$, i la que passa pels punts $B$, $C = (3, 3)$ i $D$ és $y = g(x)$, amb $g(x) = -\left(\dfrac{x-1}{2}\right)^{2} + 4$.

\subsection*{\textit{a)} \normalfont Calculeu les coordenades dels punts $A$, $B$ i $D$. \hspace*{\fill}\textnormal{[0,75~punts]}}

\pista{Cada punt és sobre un eix: imposa la coordenada que val zero.
\begin{itemize}
\item \textit{$A$} és a l'eix d'abscisses: resol $f(x) = 0$. Fixa't que $x = 0$ és una arrel; treu factor comú i resol l'equació de segon grau que queda. Tria l'arrel que correspon a la posició d'$A$ a la figura.
\item \textit{$B$} és a l'eix d'ordenades: substitueix $x = 0$ a $g(x)$.
\item \textit{$D$} és a l'eix d'abscisses: resol $g(x) = 0$ i, de les dues solucions, tria la que correspon a la posició de $D$ a la figura.
\end{itemize}

\medskip
\textbf{Atenció:} la resposta final són els \emph{tres} punts, cadascun amb les dues coordenades.}

\subsection*{\textit{b)} \normalfont Calculeu l'àrea de la zona puntejada. \hspace*{\fill}\textnormal{[1,25~punts]}}

\pista{La zona puntejada és la compresa entre les corbes $f$ i $g$, des de l'eix d'ordenades ($x = 0$) fins al punt de tall $C$ ($x = 3$). A partir de $C$ ja no és puntejada, sinó ratllada.

\medskip
L'àrea és el valor absolut de la integral de la diferència:
\[
\text{Àrea} = \left| \int_{0}^{3} \bigl(g(x) - f(x)\bigr)\,dx \right|
\]
Entre $x = 0$ i $x = 3$ les dues corbes no es tallen, així que n'hi ha prou amb una sola integral. Que $f$ toqui l'eix d'abscisses al punt $A$ no hi influeix.}

\subsection*{\textit{c)} \normalfont Els alumnes volen pintar la part puntejada de color blau i la part ratllada de color vermell. Sabent que l'àrea total del logotip és $\dfrac{175}{12}\,\text{m}^{2}$, de quin color necessitaran més pintura? \hspace*{\fill}\textnormal{[0,5~punts]}}

\pista{Aquí no cal cap integral nova: l'àrea total del logotip és la puntejada més la ratllada.
\begin{itemize}
\item Tens l'àrea total i l'àrea puntejada de l'apartat \textit{b)}: obtén la ratllada per resta.
\item Compara les dues àrees.
\end{itemize}

\medskip
\textbf{Atenció:} cal justificar la resposta: no et limitis a dir «el color és aquest».}

\end{document}
